% ============================================================
% SEÇÃO 8 – PROCEDIMENTOS DE BACKUP, SEGURANÇA E CONTINGÊNCIA
% ============================================================

\section{PROCEDIMENTOS DE BACKUP, SEGURANÇA E CONTINGÊNCIA}

O aplicativo será disponibilizado ao público por meio das plataformas oficiais de
distribuição de aplicativos móveis, Apple App Store (iOS) e Google Play Store
(Android). Ambas as lojas utilizam processos próprios e padronizados de validação,
auditoria e controle de qualidade, garantindo que o aplicativo seja distribuído com
segurança e que somente versões oficialmente assinadas estejam disponíveis para
os usuários.

O backend da aplicação utiliza como provedor de banco de dados e armazenamento
o \textbf{Supabase}, que opera sobre infraestrutura de alto desempenho com
disponibilidade global, replicação automática e mecanismos integrados de segurança
e contingência. O Supabase é executado sobre plataformas \textit{cloud}
amplamente certificadas e mantém políticas próprias de continuidade, incluindo:

\begin{itemize}
  \item Backups automáticos e criptografados, com retenção configurável;
  \item Redundância em múltiplas zonas de disponibilidade;
  \item Monitoramento contínuo da saúde dos serviços;
  \item Escalabilidade vertical e horizontal sob demanda;
  \item Conformidade com padrões reconhecidos internacionalmente, como SOC 2,
        GDPR e ISO/IEC 27001.
\end{itemize}

A segurança e integridade dos dados dos usuários são garantidas por camadas de
proteção:

\begin{itemize}
  \item Autenticação segura via JWT com tokens de acesso (curta duração) e
        refresh tokens (longa duração), implementados no backend Node.js;
  \item Renovação automática de sessão: ao receber resposta 403, o aplicativo
        solicita novo access token sem interromper o uso;
  \item Invalidação de refresh token no logout, impedindo reutilização;
  \item Criptografia em trânsito (TLS 1.2+);
  \item Criptografia em repouso (AES-256) no Supabase;
  \item Senhas armazenadas com hash + salt no banco de dados;
  \item Recuperação de senha via código de 6 dígitos enviado por e-mail: o código é armazenado apenas como hash SHA-256 com validade de 15 minutos e invalidado após uso;
  \item Moderação comunitária: sistema de denúncias com thresholds automáticos (5 denúncias ocultam uma reclamação; 10 denúncias restringem uma conta);
  \item Verificação de vínculo com a cidade: mecanismo duplo para garantir que os usuários cadastrados pertencem ao município atendido pelo aplicativo, descrito em detalhe na subseção a seguir;
  \item Políticas de acesso granular (Row-Level Security) no banco PostgreSQL;
  \item Auditorias e logs transacionais.
\end{itemize}

\subsection{Verificação de vínculo com a cidade}

Uma das premissas do LifeCity é que os dados gerados pelos usuários reflitam
problemas reais do município. Para reduzir o risco de cadastros fraudulentos ou
de usuários externos à cidade poluírem a base de dados, o sistema adota dois
mecanismos complementares:

\begin{enumerate}
  \item \textbf{Validação de CEP no cadastro (barreira de entrada):}
        Durante o processo de criação de conta, o usuário é obrigado a informar
        seu CEP. O backend valida se o CEP informado pertence à faixa de códigos
        postais do município configurado. Caso o CEP não corresponda à cidade,
        o cadastro é rejeitado com mensagem explicativa. A faixa de CEPs válidos
        é configurada por variável de ambiente (\texttt{CITY\_CEP\_PREFIXES}),
        tornando o sistema reutilizável para qualquer município. Para
        Campinas-SP, a faixa configurada abrange os prefixos 130 e 1310--1314,
        correspondentes aos CEPs 13000--13149.

  \item \textbf{Geofencing passivo por comportamento (barreira contínua):}
        A cada reclamação registrada, o backend verifica se as coordenadas
        GPS da ocorrência estão dentro dos limites geográficos do município,
        definidos por um \textit{bounding box} configurável
        (\texttt{CITY\_LAT\_MIN}, \texttt{CITY\_LAT\_MAX},
        \texttt{CITY\_LNG\_MIN}, \texttt{CITY\_LNG\_MAX}). O resultado
        (\texttt{is\_within\_city}) é armazenado junto à reclamação.
        Após o usuário acumular cinco ou mais reclamações geolocalizadas,
        o sistema verifica automaticamente o padrão: se nenhuma delas foi
        registrada dentro dos limites da cidade, a conta recebe o marcador
        \texttt{low\_trust = TRUE} no banco de dados, sinalizando ao sistema
        que aquele usuário provavelmente não pertence ao município. Essa
        verificação ocorre de forma assíncrona, sem impacto no tempo de
        resposta da criação da reclamação. Reclamações sem coordenadas GPS
        não são penalizadas.
\end{enumerate}

A combinação das duas abordagens cobre os principais vetores de abuso: o CEP
filtra no momento do cadastro com atrito mínimo para o usuário legítimo, enquanto
o geofencing passivo detecta padrões suspeitos ao longo do uso, sem exigir
verificações invasivas como documentos ou SMS. Ambos os mecanismos são
totalmente configuráveis via variáveis de ambiente, permitindo adaptar o
aplicativo a qualquer município sem alterações no código.

\bigskip

\noindent\textbf{Configuração para Campinas-SP:}
\begin{itemize}
  \item CEPs válidos: prefixos \texttt{130}, \texttt{1310}--\texttt{1314}
        (faixa 13000--13149)
  \item Limites geográficos: latitude $-23{,}15°$ a $-22{,}65°$\,S,
        longitude $-47{,}40°$ a $-46{,}90°$\,O
\end{itemize}

O plano de contingência está estruturado com base nos mecanismos nativos das
plataformas utilizadas:

\begin{itemize}
  \item Em caso de falhas no banco de dados, o Supabase executa failover
        automático e recuperação a partir de backups diários criptografados;
  \item O aplicativo é projetado para exibir mensagens de manutenção quando
        não for possível estabelecer conexão com o backend;
  \item Em caso de perda ou corrupção de dados, a restauração é realizada por
        meio dos backups automáticos e versionados do Supabase.
\end{itemize}
